\documentclass[11pt,a4paper]{article}
\usepackage[utf8]{inputenc}
\usepackage[T1]{fontenc}
\usepackage[italian]{babel}
\usepackage{hyperref}
\usepackage{enumitem}
\usepackage{csquotes}
\usepackage{geometry}
 \geometry{
 a4paper,
 total={170mm,257mm},
 left=20mm,
 top=20mm,
 }

\begin{document}


\section{Digital Forensics nel Processo Penale}

L'innesco dell'intero procedimento è la \textbf{notitia criminis} (denuncia, querela, referto ecc.), che costituisce il presupposto logico e giuridico per avviare le indagini. In informatica forense, la notitia criminis stessa contiene spesso i primi elementi tecnici, come log di sistema e segnalazioni dei provider.
Le indagini digitali devono rispettare rigorosamente procedure codificate, pena l'inutilizzabilità o la nullità della prova.

\subsection{Le Quattro Fasi del Procedimento Penale}

Dal momento in cui l'autorità giudiziaria riceve la notizia di reato fino alla sentenza definitiva, il procedimento si articola in quattro fasi principali. È importante distinguere tra ``procedimento'' (sempre presente e contenitore di tutte le fasi) e ``processo'' (fase eventuale che si apre solo con l'esercizio dell'azione penale).

\subsubsection{Fase I --- Indagini Preliminari}
Hanno la funzione di raccogliere elementi di prova sufficienti a sostenere (o escludere) l'imputazione. 
\begin{itemize}
    \item Il \textbf{PM} dirige l'attività della PG e nomina i Consulenti Tecnici (CTPM) per valutare se procedere o archiviare.
    \item Se il PM richiede l'archiviazione, la parte offesa può opporsi (spesso grazie all'ausilio di un proprio CTP che evidenzia lacune investigative).
    \item Se gli elementi raccolti sono sufficienti, il PM esercita l'azione penale chiedendo il rinvio a giudizio.
\end{itemize}

\subsubsection{Fase II --- Udienza Preliminare (Fase eventuale)}
Funziona da filtro processuale tra indagini e dibattimento per i reati di competenza del Tribunale collegiale.
\begin{itemize}
    \item Il Giudice dell'Udienza Preliminare (GUP) valuta, in un contraddittorio prevalentemente cartolare sugli atti d'indagine, se rinviare a giudizio l'imputato o pronunciare una sentenza di non luogo a procedere.
    \item In questa fase, l'imputato può chiedere riti alternativi (es. giudizio abbreviato o patteggiamento). Per i reati minori (citazione diretta a giudizio), questa fase non è prevista.
\end{itemize}

\subsubsection{Fase III --- Atti Preliminari al Dibattimento}
Costituisce il segmento di atti preparatori tra il decreto che dispone il giudizio e l'apertura effettiva del dibattimento.
\begin{itemize}
    \item Si formalizzano le ultime regolarità e la costituzione in giudizio delle parti (parte civile, responsabile civile).
    \item \textbf{Istanze probatorie:} Le parti depositano e formulano al Giudice le proprie richieste di prova, tra cui liste dei testimoni, la produzione documentale o le formali richieste di perizia tecnica. 
\end{itemize}

\subsubsection{Fase IV --- Il Dibattimento e la Decisione}
Il dibattimento è il cuore del processo: qui ha luogo la formazione delle prove in pieno contraddittorio tra le parti.
\begin{itemize}
    \item \textbf{Istruttoria dibattimentale:} I mezzi di prova (come l'esame di testimoni, periti e consulenti) vengono assunti davanti al Giudice in udienza.
    \item \textbf{Poteri del Giudice (Art. 507 c.p.p.):} Il Giudice ha il potere di disporre l'assunzione di ulteriori prove d'ufficio se risultano assolutamente necessarie per completare l'accertamento.
    \item \textbf{Sentenza:} Conclusa l'istruttoria e dopo la requisitoria del PM e le difese degli avvocati, il Giudice si ritira in camera di consiglio e pronuncia la sentenza, corredata da una motivazione tecnica e giuridica sui punti controversi. Se il Giudice dovesse discostarsi dal parere tecnico del Perito o del CT, è tenuto a motivare adeguatamente la sua scelta.
\end{itemize}

\subsection{Attori e Ausiliari Tecnici}
\begin{itemize}
    \item \textbf{PG (Polizia Giudiziaria)}: Interviene per prima sulla scena (es. Polizia Postale). Il suo compito primario è il ``triage'' forense: isolare il dispositivo, documentare fotograficamente l'ambiente, e ``congelare'' gli elementi. 
    \item \textbf{PM (Pubblico Ministero)}: Dominus delle indagini. Coordina la PG e valuta gli elementi raccolti per decidere sull'azione penale.
    \item \textbf{CTPM (Consulente Tecnico del PM)}: Nominato dal PM per affiancare la polizia. Lavora per l'accusa con l'obbligo della verità scientifica.
    \item \textbf{Difesa e CTP (Consulente Tecnico di Parte)}: Controlla e confuta le tesi dell'accusa. Verifica il rispetto della legge sulle metodologie di acquisizione e può assistere alle operazioni del CTPM, presentare memorie tecniche e formulare riserve se nota difetti nella catena di custodia.
    \item \textbf{Giudice e Perito}: Il Perito è un esperto super-partes che interviene solitamente in fase dibattimentale per dipanare questioni tecniche complesse. Risponde ai quesiti del Giudice e il suo elaborato fa piena prova.
\end{itemize}

\subsection{Tipologie di Accertamento Tecnico}
La distinzione tra ripetibile e irripetibile è il cuore della procedura penale informatica:
\begin{itemize}
    \item \textbf{Accertamenti Ripetibili}: Sono eseguiti sulle copie forensi. Poiché l'originale non viene alterato e l'operazione può essere rifatta identica mille volte, il PM procede senza preavvisare la difesa.
    \item \textbf{Accertamenti Irripetibili}: Interventi che alterano il dato o l'hardware. Impongono l'obbligo di preavviso a indagati e parte offesa.
    \item \textbf{Incidente Probatorio}: Una ``parentesi'' delle indagini preliminari. Se c'è rischio che una prova si disperda o modifichi prima del processo, viene assunta subito davanti al Giudice (GIP) con tutte le garanzie del dibattimento.
\end{itemize}

\subsection{Mezzi di Ricerca della Prova Digitale}
\begin{itemize}
    \item \textbf{Ispezione}: Attività descrittiva su un sistema per rilevarne le tracce e gli effetti materiali. Spesso fatta su sistemi ``live''.
    \item \textbf{Perquisizione}: Ricerca mirata di dati o dispositivi specifici. È un atto a sorpresa.
    \item \textbf{Sequestro}: Acquisizione coatta del dispositivo o del dato. Nel digitale, sempre più spesso si sequestra il ``dato'' (es. tramite cloud extraction o acquisizione presso il provider) restituendo l'hardware per minimizzare il danno all'azienda.
    \item \textbf{Intercettazioni Telematiche}: Captazione occulta del flusso dati in transito (es. trojan di stato, sniffing di rete).
\end{itemize}

\subsection{La Relazione Tecnica}
Il risultato di qualsiasi accertamento deve confluire in una relazione tecnica.
\begin{itemize}
    \item Deve essere \textbf{comprensibile}: Il Giudice è ``peritus peritorum'' ma non è un informatico.
    \item Deve indicare la \textbf{metodologia}: Strumenti usati, versioni software, algoritmi di hashing.
    \item Deve distinguere tra \textbf{fatto tecnico} (``il file X è stato creato il giorno Y'') e \textbf{valutazione} (``è probabile che l'utente sapesse'').
\end{itemize}

\subsection{In Dibattimento}
L'esame incrociato (cross-examination) è il momento della verità. 
Il CTP cercherà di invalidare la prova attaccando la metodologia, la catena di custodia o le conclusioni. 
Un esperto di parte ben preparato può far emergere errori o imprecisioni, rendendo la prova inutilizzabile.

% --------------------------------------------------

\section{Digital Forensics nel Processo Civile}

A differenza del processo penale, il processo civile si occupa di controversie tra soggetti privati (persone fisiche o giuridiche) per la tutela di un diritto o il risarcimento di un danno, e non a punire un reato. 
Gli ambiti tipici includono concorrenza sleale, violazione di patti di non concorrenza, furto di proprietà intellettuale, diffamazione, o inadempienze nei contratti IT.
L'obiettivo della Digital Forensics in questo contesto è fornire evidenze oggettive, verificabili e inalterabili a supporto delle pretese economiche o legali delle parti.

\subsection{Princìpi Fondamentali della Prova}

\subsubsection{L'Onere della Prova}
Nel processo civile vige il principio cardine secondo cui \textit{``Chi vuol far valere un diritto in giudizio deve provare i fatti che ne costituiscono il fondamento''}. Questo significa che spetta all'Attore portare al Giudice prove informatiche tecnicamente inattaccabili.

\subsubsection{La Prova di Resistenza}
Poiché il dato digitale è intrinsecamente volatile e facilmente manipolabile, l'autenticità della prova va accertata oltre ogni ragionevole dubbio, altrimenti non ha alcun valore probatorio se disconosciuta formalmente dalla controparte.
La prova deve superare una ``prova di resistenza'' dimostrando la sua inalterabilità tramite firme digitali, marche temporali, estrazione degli header originali, calcolo dell'hash e catena di custodia certificata.

\subsection{Attori e Ausiliari Tecnici}

\begin{itemize}
    \item \textbf{Attore e Convenuto}: Le parti private del processo. L'Attore avvia la causa citando in giudizio la controparte, il Convenuto si costituisce per difendersi.
    \item \textbf{Giudice}: Dirige il processo, valuta le memorie, ammette le prove e prende la decisione finale.
    \item \textbf{CTU (Consulente Tecnico d'Ufficio)}: Previsto dall'Art. 61 c.p.c., è l'esperto nominato dal Giudice quando per la decisione sono necessarie specifiche competenze tecniche. A differenza del perito in ambito penale, non è un autonomo mezzo di prova, ma uno strumento per integrare le conoscenze del Giudice. Risponde rigorosamente ed esclusivamente ai ``quesiti'' formulati dal magistrato e deve operare in totale imparzialità.
    \item \textbf{CTP (Consulenti Tecnici di Parte)}: Nominati da Attore e Convenuto. Affiancano il CTU in tutte le operazioni, tutelano gli interessi dei propri clienti, verificano la correttezza metodologica delle estrazioni e presentano osservazioni tecniche vincolanti alla bozza di relazione del CTU.
\end{itemize}

\subsection{Strumenti Pre-Contenziosi e d'Urgenza}

Spesso nel contenzioso civile la prova digitale rischia di essere distrutta o alterata prima che la causa formale di merito entri nel vivo. Per questo l'ordinamento prevede strumenti specifici:
\begin{itemize}
    \item \textbf{Accertamento Tecnico Preventivo}: Viene nominato un CTU per ``congelare'' la situazione, estrarre i dati e quantificare danni prim'ancora di avviare il processo vero e proprio (usato tipicamente nei contenziosi per software malfunzionante). L'esito oggettivo dell'ATP porta molto spesso le parti a chiudere la disputa con una transazione stragiudiziale bonaria.
    \item \textbf{Descrizione}: Una sorta di ``blitz forense''. Se un'azienda sospetta il furto di segreti commerciali da parte di un concorrente, chiede al Giudice l'autorizzazione a ispezionarne i sistemi aziendali \textit{inaudita altera parte} (senza alcun preavviso). Un Ufficiale Giudiziario, affiancato dal CTU, si presenta a sorpresa per eseguire l'acquisizione forense mirata dei dati.
    \item \textbf{Ricorso d'Urgenza}: Azione intrapresa quando vi è un \textit{periculum in mora} (pericolo grave e imminente), ad esempio se c'è il sospetto fondato che un dipendente dimissionario stia attualmente formattando i database aziendali per portarli altrove.
\end{itemize}

\subsection{Le Fasi dell'Accertamento Tecnico (CTU)}

Quando il Giudice dispone una CTU informatica, il percorso si articola in fasi procedurali garantiste del pieno contraddittorio:

\subsubsection{1. Nomina e Formulazione del Quesito}
Il Giudice nomina il CTU, il quale presta giuramento di rito. Nella stessa udienza, il Giudice definisce formalmente il ``Quesito'' (es. ``Accerti il CTU se dai computer in uso al convenuto vi sia stata esfiltrazione di dati verso memorie USB''). Le parti, di conseguenza, nominano i propri CTP.

\subsubsection{2. Inizio Operazioni Peritali (IOP) e Acquisizione}
È il momento in cui avvengono le acquisizioni delle copie forensi. A differenza del penale (spesso caratterizzato da sequestri a sorpresa), nel civile tutto deve svolgersi in contraddittorio palese. Il CTU convoca ufficialmente i CTP e l'acquisizione dei dispositivi avviene alla presenza di tutti. Si adottano \textit{write-blocker}, si isolano le evidenze e si calcolano gli algoritmi di hash per cristallizzare la prova e impedire contestazioni future sull'integrità del reperto.

\subsubsection{3. Analisi, Espunzione e Limiti della Privacy (GDPR)}
Questa fase segna una demarcazione fortissima rispetto al penale. Nel civile \textbf{non sono ammesse} ispezioni massive e indiscriminate sui dispositivi. 
Il CTU è tenuto ad applicare rigidi filtri di pertinenza e a \textbf{espungere} tutte le email personali, i dati medici o i file irrilevanti per il contenzioso. 
Tutto ciò nel rispetto assoluto del GDPR. Violazioni della privacy in questa fase rendono nulla l'intera prova ed espongono l'esperto a pesanti sanzioni.

\subsubsection{4. Bozza e Osservazioni}
Conclusa l'analisi tecnica, il CTU redige una bozza preliminare di perizia e la invia esclusivamente ai CTP. Questi ultimi hanno un termine stabilito per rispondere con le loro ``Osservazioni Tecniche'', mediante le quali possono attaccare le metodologie usate dal CTU, evidenziare lacune o proporre interpretazioni alternative a favore dei propri assistiti.

\subsubsection{5. Relazione Finale}
Il CTU deposita in Tribunale la relazione definitiva, che deve obbligatoriamente contenere e contro-dedurre in modo puntuale alle osservazioni sollevate dai CTP. Questa relazione diventa un pilastro istruttorio decisivo su cui il Giudice (che pur essendo \textit{peritus peritorum} non possiede le competenze informatiche) fonderà la propria sentenza.

% --------------------------------------------------

\section{Differenze tra Processo Penale e Processo Civile in Digital Forensics}
\subsection{1. Scopo dell'Indagine e Iniziativa}
\begin{itemize}
    \item \textbf{Processo Penale}: L'obiettivo è l'accertamento della verità materiale per verificare se è stato commesso un reato e punire il colpevole. L'azione penale è \textbf{obbligatoria} e viene esercitata dallo Stato tramite il Pubblico Ministero (PM) non appena vi è una \textit{notitia criminis}.
    \item \textbf{Processo Civile}: L'obiettivo è dirimere una controversia patrimoniale o tutelare un diritto tra soggetti privati. Il Giudice non cerca la verità assoluta, ma decide in base alle prove fornite dalle parti. L'iniziativa è \textbf{di parte}: il processo inizia solo se un soggetto (l'Attore) decide volontariamente di citare in giudizio un altro (il Convenuto).
\end{itemize}

\subsection{2. L'Onere della Prova}
\begin{itemize}
    \item \textbf{Processo Penale}: Spetta all'accusa (il PM) dimostrare la colpevolezza dell'imputato \textbf{``oltre ogni ragionevole dubbio''}. Se la prova digitale (es. l'attribuzione di un indirizzo IP a una persona fisica) presenta zone d'ombra o difetti di catena di custodia, il Giudice deve propendere per l'assoluzione.
    \item \textbf{Processo Civile}: Vige l'Art. 2697 c.c.: \textbf{``Chi vuol far valere un diritto in giudizio deve provare i fatti che ne costituiscono il fondamento''}. L'Attore deve produrre prove digitali inoppugnabili (superando la ``prova di resistenza'') e il Giudice decide basandosi sulla regola del ``più probabile che non'' (preponderanza dell'evidenza).
\end{itemize}

\subsection{3. Gli Ausiliari Tecnici del Giudice}
\begin{itemize}
    \item \textbf{Il Perito (Penale)}: Viene nominato dal Giudice eccezionalmente, solo quando occorre svolgere indagini che richiedono specifiche competenze tecniche che esulano dal patrimonio conoscitivo del magistrato. Il suo elaborato costituisce un \textbf{mezzo di prova autonomo}.
    \item \textbf{Il CTU (Civile)}: Viene nominato dal Giudice quasi sistematicamente nelle cause tecnologiche. Non è un mezzo di prova, ma uno \textbf{strumento istruttorio ed estensione tecnica del Giudice}. Il CTU non può indagare liberamente, ma è strettamente vincolato al ``quesito'' formulato dal Giudice.
\end{itemize}

\subsection{4. Modalità di Acquisizione e Contraddittorio}
\begin{itemize}
    \item \textbf{In Ambito Penale}: L'acquisizione della prova è spesso \textbf{coatta e a sorpresa} (ispezione, perquisizione, sequestro tramite Polizia Giudiziaria). Il contraddittorio può essere posticipato: nel caso di accertamenti irripetibili la difesa viene avvisata, ma nel caso di accertamenti ripetibili il PM procede senza avvisare la controparte.
    \item \textbf{In Ambito Civile}: Domina il \textbf{contraddittorio preventivo e continuo}. Le acquisizioni forensi avvengono durante l'Inizio Operazioni Peritali (IOP) indetto dal CTU, alla presenza dei CTP di entrambe le parti, che possono vigilare in tempo reale su ogni singolo click. Le uniche eccezioni a sorpresa sono gli ordini di ``Descrizione'' \textit{inaudita altera parte} per tutelare i segreti industriali.
\end{itemize}

\subsection{5. Limiti dell'Indagine, Pertinenza e Privacy}
\begin{itemize}
    \item \textbf{L'Indagine Penale}: Tende a essere ``massiva''. Se viene sequestrato il computer di un indagato, il CTPM clona l'intero hard disk e analizza tutto il contenuto alla ricerca di qualsiasi traccia pertinente al reato, con ampi poteri di ispezione.
    \item \textbf{L'Indagine Civile}: È dominata dal principio di stretta pertinenza e dal GDPR. Sono assolutamente vietate le \textit{fishing expedition} (acquisizioni a strascico sperando di trovare qualcosa di illecito). Il CTU è obbligato a \textbf{espungere} (isolare e non acquisire/cancellare) tutto ciò che attiene alla sfera personale, sanitaria o estranea al mandato, pena la nullità della prova.
\end{itemize}

% --------------------------------------------------

\section{Firme Elettroniche}
La gestione dei documenti informatici si basa su un impianto normativo integrato tra l'\textbf{eIDAS} e il \textbf{CAD} (Codice dell'Amministrazione Digitale).

Ogni accertamento su firme e documenti digitali poggia su tre proprietà fondamentali:
\begin{itemize}
    \item \textbf{Integrità}: corrispondenza esatta tra originale ed esibito, verificata tramite hashing.
    \item \textbf{Immodificabilità}: garanzia che l'informazione non sia stata alterata nel tempo.
    \item \textbf{Paternità}: connessione univoca tra il bitstream e l'identità del firmatario tramite chiavi crittografiche (basate su PKI).
\end{itemize}

\subsection{Classificazione delle Firme Elettroniche}
Il regolamento eIDAS definisce una gerarchia a 3 livelli:
\begin{itemize}
    \item \textbf{FES (Firma Elettronica Semplice)}: dati in formato elettronico associati a un documento (es. PIN, nome in email). È liberamente valutabile dal giudice e possiede valore probatorio limitato.
    \item \textbf{FEA (Firma Elettronica Avanzata)}: possiede 4 requisiti (art. 26 eIDAS), è connessa univocamente al firmatario, idonea a identificarlo e creata con dati sotto il suo controllo esclusivo. Ha il valore di scrittura privata.
    \item \textbf{FEQ (Firma Elettronica Qualificata)}: FEA + certificato qualificato + QSCD. Prevede un'identificazione fisica al momento del rilascio e si appoggia a terzi come garanzia.
\end{itemize}

\subsubsection{Standard italiano} 
Sottospecie della FEQ basata su infrastruttura PKI nazionale e crittografia asimmetrica (RSA/ECDSA). La firma si applica sull'hash del documento. 
Gode dell'\textbf{inversione dell'onere della prova}: per disconoscerla, il titolare deve allegare fatti specifici e dimostrare l'uso abusivo del dispositivo.

\subsubsection{Formati di Firma}
\begin{itemize}
    \item \textbf{PAdES (PDF)}: firma incorporata nel PDF, ideale per contratti.
    \item \textbf{CAdES (.p7m)}: busta CMS che ingloba documento e firma in un unico file. È lo standard de facto nel B2B e nella PA italiana.
    \item \textbf{XAdES (XML)}: firma su strutture XML (es. FatturaPA) per flussi interoperabili.
\end{itemize}

\subsection{Validazione e Marcatura Temporale}
La validazione di una firma si articola in: 
\begin{enumerate}
    \item Validità del certificato al time-of-signing
    \item Verifica del path di certificazione
    \item Verifica crittografica
\end{enumerate}

\begin{quote}
\textbf{Marcatura Temporale (Time Stamping):}  Attestazione di una Timestamp Authority (TSA) che garantisce che un documento esisteva in una certa forma ad un dato istante. La marca è legata crittograficamente all'hash.
\end{quote}

Garantisce la possibilità di verificare la firma al momento della sottoscrizione, anche a certificato scaduto, risultando essenziale per la Long Term Validation (LTV).

\subsection{Analisi delle Copie}
Il CAD distingue rigidamente tra le repliche di un documento informatico:

\begin{itemize}
    \item \textbf{Duplicato Informatico (1:1)}: Ha la stessa identità assoluta a livello di bitstream e il medesimo valore giuridico dell'originale.
    \item \textbf{Copia Informatica}: contenuto identico ma sequenza binaria diversa (es. cambio formato, aggiornamento metadati di sistema). Ha la stessa efficacia dell'originale solo se attestata da un pubblico ufficiale o se non espressamente disconosciuta.
    \item \textbf{Copia Analogica di Doc. Informatico (Rematerializzazione)}: la stampa perde irrimediabilmente la firma crittografica. L'autenticità è recuperabile solo tramite un \textbf{contrassegno} (QR Code o glifo) che collega la copia cartacea all'originale digitale.
\end{itemize}

\subsection{Prassi e Analisi Forense}
\subsubsection{FES ed Email}
Essendo la FES liberamente valutabile, il Consulente Tecnico (CT) deve fornire criteri oggettivi di affidabilità attraverso:
\begin{itemize}
    \item Tracciabilità dell'accesso (log, IP).
    \item Coerenza temporale (artefatti di sistema).
    \item Identificazione del soggetto (pattern comportamentali).
\end{itemize}

Per l'email (considerata documento informatico), l'analisi forense prevede l'esame sistematico di:
\begin{enumerate}
    \item \textbf{Header RFC 5322}: per analizzare campi \textit{Received} e anomalie degli IP (spoofing).
    \item \textbf{Protocolli (SPF, DKIM, DMARC)}: il loro allineamento rappresenta la prova più forte di autenticità.
    \item \textbf{Struttura MIME}: per rilevare alterazioni post-invio tramite l'analisi dei boundary.
\end{enumerate}

\subsubsection{PDF Incremental Updates}
L'architettura dei PDF firmati prevede che non si sovrascrivano i dati. Le modifiche avvengono per sezioni incrementali: la firma copre solo i byte precedenti all'apposizione, permettendo al perito di estrarre e visualizzare le versioni esatte pre-modifica tramite byte-offset.

\subsection{Conservazione Digitale}
Obbligatoria per la PA e facoltativa per i privati (ruoli specifici previsti), serve a mantenere il valore probatorio nel tempo. Poiché i certificati scadono e gli algoritmi crittografici diventano obsoleti, il conservatore qualificato utilizza il chained timestamping e i formati LTA (Long Term Archival) applicati al \textbf{Pacchetto di Archiviazione (PdA)}.

% --------------------------------------------------

\section{Decreto legislativo 231/2001}
Legge cambia il trattamento verso organizzazioni, abbandonando il concetto della colpa attribuita esclusivamente alla persona fisica, ed estendendo le responsabilità ad organizzazioni che beneficiano delle condotte illecite dei singoli componenti.
Inizialmente applicata esclusivamente alla corruzione e successivamente espansa. Alcuni esempi rilevanti sono: sicurezza sul lavoro, utilizzo abusivo di sistemi informatici...

L'integrazione dei reati informatici ha reso necessaria una competenza tecnica specialistica nei modelli di compliance. Il sistema sanzionatorio prevede:
\begin{itemize}
    \item \textbf{Sanzioni pecuniarie a quote}: commisurate alla capacità economica dell'ente (da un minimo di 100 quote a oltre 1,5 milioni di euro per i casi più gravi).
    \item \textbf{Sanzioni interdittive}: possono paralizzare l'attività (esclusione da finanziamenti, divieto di contrattare con la PA, sospensione licenze).
\end{itemize}

\subsection{Modello Organizzativo - MOG}
Le aziende hanno l'obbligo di istituire \textbf{Modelli Organizzativi (MOG)} per prevenire questo tipo di reato da parte dei componenti. Esso è costituito da tutti gli accorgimenti tecnici/organizzativi atti a prevenire condotte illecite.
Dal punto di vista informatico si istituisce quindi l'obbligo in materie come:
\begin{itemize}
    \item Monitoraggio degli impianti e asset inventory.
    \item Certificazione digitale delle attività di formazione.
    \item Documentare la contabilità e i flussi finanziari.
    \item Protocolli tecnici: Segregazione delle funzioni (SoD), Identity \& Access Management (IAM), Crittografia e MFA.
\end{itemize}

\begin{quote}
In un contesto processuale quanto queste misure siano sufficienti è appannaggio del giudice (affiancato dai CT). La validità del Modello non si misura sulla sua esistenza formale, ma sulla sua idoneità concreta e verificabile: serve la prova tecnica dell'attuazione (log, audit, penetration test).
\end{quote}

In generale la ratio è dimostrare di aver preso ogni misura possibile per prevenire reati da parte dei dipendenti (o danni da loro subiti).

Inoltre si istituisce l'obbligo da parte delle organizzazioni di costituire un \textbf{Organismo di Vigilanza (ODV)} dotata di autonomi poteri di iniziativa e controllo.

In maniera periodica o in seguito ad alert/whistleblowing l'ODV ha l'obbligo di condurre attività di auditing. 
\begin{quote}
L'attività di auditing spesso include anche l'analisi forense di singoli dispositivi dei dipendenti selezionati randomicamente. Per l'efficacia del MOG digitale è fondamentale la Log Retention su supporti WORM (Write Once Read Many) per garantire l'integrità dei dati.
\end{quote}

\section{GDPR e gestione di data breach}
Definizioni:
\begin{itemize}
    \item \textbf{Dato personale}: qualsiasi informazione relativa a una persona identificabile.
    \item \textbf{Trattamento dei dati}: qualsiasi operazione sui dati. \textbf{NB}: anche l'archiviazione è trattamento.
    \item \textbf{Pseudoanonimizzazione}: trasformazione dei dati atta a rendere più difficile l'associazione del dato alla persona fisica.
\end{itemize}

\subsection{Figure e responsabilità}
\begin{itemize}
    \item \textbf{Titolare del trattamento}: Stabilisce le modalità e mezzi del trattamento. È responsabile verso il garante e gli interessati e ha l'obbligo di dimostrare l'adempimento delle misure con evidenze documentali.
    \item \textbf{Responsabile del trattamento}: Soggetto che tratta dati per conto del titolare, garantendo misure tecniche adeguate.
    \item \textbf{Amministratore di sistema}: Figura tecnica con privilegi elevati, critica per la sicurezza dei log e degli accessi.
    \item \textbf{Data Protection Officer (DPO)}: ha funzione di supporto e si interfaccia col garante. Collabora con CISO e ODV per far convergere la compliance 231 con la protezione dati.
\end{itemize}

\subsection{Sanzioni}
Le sanzioni sono calibrate sulla gravità della violazione.
\begin{itemize}
    \item Art. 83 par. 4 - Non adempimento di obblighi tecnico-organizzativi: fino a 10M o 2\% fatturato globale.
    \item Art. 83 par. 5 - Violazione dei principi di base, diritti degli interessati, trasferimenti, ordini del Garante: fino a 20M o 4\% fatturato globale.
    \item Danno agli interessati - Responsabilità solidale titolare-responsabile: Risarcimento civile (Art. 82).
\end{itemize}

Gli aggravanti/attenuanti dell'84 par. 2 includono la cooperazione con l'autorità e l'adozione di misure adottate per mitigare il danno.

\subsubsection{Privacy by Design \& Default}
La protezione dei dati deve essere tenuta in considerazione in tutte le fasi di della progettazione, implementazione e deployment di un sistema informatico (Art. 25 GDPR).

\subsubsection{Diritti dell'interessato}
L'interessato ha diritto ad:
\begin{itemize}
    \item \textbf{Accesso}: copia dei propri dati.
    \item \textbf{Oblio}: cancellazione sicura e definitiva.
    \item \textbf{Portabilità}: trasferimento dei dati in formato strutturato e leggibile da dispositivo automatico.
\end{itemize}

\subsection{Data Protection Impact Assessment - DPIA}
Il DPIA è uno strumento \textbf{ex ante obbligatorio} per trattamenti che presentano criticità per diritti/libertà delle persone fisiche. Corrisponde a un'analisi sistematica dei rischi per identificare le misure tecniche necessarie a mitigarli prima dell'inizio del trattamento.

\section{Data Breach}
Violazione della sicurezza che comporta la perdita di: \textbf{Confidenzialità, Integrità o Disponibilità}.

La notifica al Garante è evitabile se il titolare valuta (documentando) che il rischio per i diritti e le libertà delle persone fisiche è improbabile.

Altrimenti è obbligo la comunicazione al Garante entro 72h dalla scoperta dell'incidente. Va documentato:
\begin{itemize}
    \item Tipologia di dati e contenuto.
    \item Numero e categoria degli interessati.
    \item Misure di mitigazione intraprese.
\end{itemize}

\subsection{NIS2 e DORA}
Normative europee che impongono obblighi di cybersecurity per servizi essenziali e settore finanziario. Richiedono:
\begin{itemize}
    \item Incident reporting (notifica entro 24h).
    \item Gestione del rischio della supply chain.
    \item Test di resilienza operativa.
\end{itemize}
La loro violazione costituisce elemento di ``colpa organizzativa'' rilevante ai fini della responsabilità 231.

\subsection{Procedure operative e Investigazioni Interne}
A seguito di un alert o whistleblowing, l'azienda deve cristallizzare immediatamente la prova digitale.
\begin{itemize}
    \item \textbf{Paziente Zero}: analisi per ricostruire il vettore d'attacco.
    \item \textbf{Chain of Custody}: fondamentale per la difesa processuale. Si seguono standard come RFC 3227 (acquisizione RAM, immagini forensi bit-by-bit, verifica hash).
    \item \textbf{Finalità}: dimostrare l'estraneità dell'ente o l'elusione fraudolenta del MOG da parte dell'autore.
    \item \textbf{Whistleblowing}: obbligo di piattaforme cifrate (es. GlobaLeaks) per garantire la riservatezza assoluta del segnalante.
\end{itemize}

\subsection{Cyber Resilience Act (CRA)}

Il \textbf{Cyber Resilience Act (CRA)} è un regolamento europeo che introduce requisiti obbligatori di sicurezza informatica per i ``prodotti con elementi digitali'' (hardware e software) immessi sul mercato dell'Unione Europea. Si affianca ad altre normative strategiche (come la Direttiva NIS2 e il DORA) per rafforzare la postura di sicurezza complessiva del mercato unico.

L'obiettivo normativo primario è spostare il peso e la responsabilità della sicurezza informatica dagli utenti finali (consumatori o aziende) ai \textbf{produttori}, garantendo che i dispositivi e i software siano protetti fin dalla fabbrica e per tutto il loro ciclo di vita.

\subsubsection{Ambito di Applicazione}
La normativa ha una portata vastissima: si applica a tutti i prodotti hardware e software connessi, direttamente o indirettamente, a una rete o a un altro dispositivo (es. dispositivi IoT, router, telecamere, smart appliance, sistemi operativi, software gestionale). 
\begin{quote}
\textbf{NB:} Sono esentati dal CRA i settori che sono già ampiamente regolamentati da normative specifiche sulla sicurezza, come i dispositivi medici, il settore automotive e l'aviazione civile.
\end{quote}

\subsubsection{Obblighi Fondamentali dei Produttori}
Il CRA impone regole stringenti che trasformano le best practice tecniche in obblighi di legge:
\begin{itemize}
    \item \textbf{Security by Design \& Default:} I prodotti devono essere progettati per mitigare i rischi e distribuiti con le impostazioni di massima sicurezza attive di default (senza richiedere l'intervento dell'utente).
    \item \textbf{Assenza di vulnerabilità note:} Al momento dell'immissione sul mercato, i prodotti non devono presentare vulnerabilità già note e sfruttabili.
    \item \textbf{Gestione e Aggiornamenti:} È fatto obbligo di fornire aggiornamenti di sicurezza gratuiti, tempestivi e (ove possibile) automatici per un periodo adeguato all'uso previsto del prodotto, tipicamente fissato fino a 5 anni o per il ciclo di vita atteso.
\end{itemize}

\subsubsection{Implicazioni per la Digital Forensics e la Compliance}
L'introduzione del CRA ha un impatto profondo sulle procedure di Digital Forensics e sulla dimostrazione della compliance (es. Modelli 231 e GDPR), fornendo agli investigatori nuovi strumenti e artefatti documentali:
\begin{itemize}
    \item \textbf{SBOM (Software Bill of Materials):} I produttori devono creare e mantenere una ``distinta base'' del software, mappando tutte le dipendenze e le librerie (incluse quelle open-source). In ambito forense, la SBOM è cruciale: permette a un CTU o a un analista di verificare immediatamente se un dispositivo compromesso conteneva librerie vulnerabili note che il produttore ha trascurato di aggiornare.
    \item \textbf{Incident Reporting Accelerato:} Esiste l'obbligo di notificare all'ENISA (Agenzia dell'UE per la cybersicurezza) e ai CSIRT nazionali qualsiasi vulnerabilità attivamente sfruttata o incidente grave entro \textbf{24 ore} dalla scoperta. Questo restringe drasticamente i tempi operativi rispetto alle 72 ore previste dal GDPR.
    \item \textbf{Tracciabilità e Log di Sicurezza:} I dispositivi devono essere ingegnerizzati per proteggere l'integrità dei dati e fornire meccanismi di registrazione (log retention) adeguati per consentire il monitoraggio e la successiva analisi forense post-incidente.
\end{itemize}

\subsection{Quadro Sanzionatorio}
Così come il GDPR e il D.Lgs 231/01, il CRA utilizza la leva economica per imporre l'adeguamento. Le violazioni dei requisiti essenziali di sicurezza prevedono un apparato sanzionatorio molto severo, che può arrivare fino a \textbf{15 milioni di euro} oppure al \textbf{2,5\% del fatturato globale annuo} dell'azienda (se tale importo risulta superiore).

% --------------------------------------------------

\section{Privacy del lavoratore}

\begin{quote}
\textbf{Extra.} Non è possibile redigere accordi scritti tra datore di lavoro e dipendente in quanto viene data per scontata la condizione di sottomissione/soggezione tra le parti.
\end{quote}

Il controllo automatizzato/a distanza dei lavoratori è generalmente proibito e costituisce un reato.
\begin{quote}
\textbf{NB.} Anche la predisposizione di sistemi di sorveglianza senza che siano utilizzati costituisce un reato (crea comunque una situazione di soggezione).
\end{quote}

Per essere consentito deve essere autorizzato preventivamente secondo doppio canale:
\begin{itemize}
    \item \textbf{Sindacale}: I sindacati rappresentano i dipendenti nella stipulazione di accordi per l'istituzione di misure di sorveglianza 
    \item \textbf{Amministrativa}: In mancanza di accordi sindacali è possibile fare un istanza all'ispettorato del lavoro (ITL/INL). Va presentato un progetto che documenti il tipo e finalità delle misure di sorveglianza.
\end{itemize}

\begin{quote}
Per la videosorveglianza c'è un limite di conservazione di 48h.
\end{quote}

Tutte le regolamentazioni in materia sono definite dallo statuto del lavoratore.

\subsection{Misure di controllo consentite}
Gli unici casi in cui è consentita la sorveglianza sono la protezione del patrimonio aziendale e la tutela della sicurezza del lavoratore.

\subsubsection{Presupposto dell'informativa}
Qualsiasi misura di sorveglianza presuppone:
\begin{itemize}
    \item \textbf{Preventiva e Adeguata informativa}: il dipendente deve essere informato in maniera specifica sulle misure di controllo e con una finestra di tempo sufficiente prima che questi vengano messi in atto.
    \item \textbf{Rispetto del GDPR}: il trattamento dei dati raccolti deve essere conforme al GDPR.
\end{itemize}

\subsubsection{Controlli difensivi}
Controlli mirati ad accertare specifiche condotte illecite già sospettate. Il sospetto deve pre-esistere al controllo e deve essere oggettivamente verificabile. Non è possibile avviare indagini forensi generiche al fine di cercare irregolarità.

L'implementazione di queste misure deve essere preceduta dall'apertura di un fascicolo d'indagine difensiva e viene curata da un avvocato.

\subsubsection{Tipologia di controlli}
\begin{itemize}
    \item \textbf{Controlli volontari}: Metodologie di controllo introdotte volontariamente in tecnologie progettate per uso generico, ie. GPS su veicoli/RFID nei badge. Passano obbligatoriamente attraverso i sindacati o ITL/INL.
    \item \textbf{Controlli involontari}: Funzioni di controllo presenti sui dispositivi ma non esplicitamente impiegate per questo scopo, ie. GPS mobile/webcam/microfoni integrati...
\end{itemize}

\begin{quote}
L'analisi del traffico di rete è consentita solamente per finalità diagnostiche. Analogamente misure analoghe al jamming sono sanzionate penalmente.
\end{quote}

\subsection{Esonero dall'autorizzazione}
Alcune casistiche non richiedono la previa autorizzazione da sindacati o ITL/INL.
\begin{itemize}
    \item Cronotachigrafo
    \item Telepass su veicolo aziendale
    \item Geolocalizzazione di smartphone aziendale
\end{itemize}

\subsubsection{Territorio grigio}
Tecnologie moderne introducono nuove difficoltà normative. Dispositivi come terminali personali/smart glove/smart watch rendono difficile (se non impossibile) distinguere l'attività di accertamento forense ed investigazioni illegittime.

% --------------------------------------------------
\section{Segreti commerciali}

\begin{quote}
Prerequisiti per l'accesso abusivo 
Presenza di  misure di sicurezza (ius escludenti, principio labile), in contesto penale è sufficiente la loro presenza, in contesto civile è necessario che queste misure siano adeguate.
\end{quote}

\subsection{Segretezza vs Riservatezza}
Nel contesto di rapporti B2B/rapporti col cliente/dipendenti vi è la possibilità di concordare NDA, obblighi di riservatezza unilaterali o reciproci che impongono il divieto di diffusione di una certa informazione. Tipicamente vengono previste allo stesso tempo obblighi e penali.

\subsection{Codice della proprietà industriale}
Canonizza i principi sulla preservazione sul patrimonio immateriale di un azienda (tecniche operative/processi aziendali/competenze di settore/...).
Per informazioni segrete commerciali (fornitori e accordi con questi/pricing models/margini/clientela/...) si intendono in generale informazioni che hanno rilevanza economica.

\subsubsection{Art. 98 CPI - Oggetto della tutela}
A livello civilistico per rientrare nella tutela dei segreti un'informazione deve rispettare tre requisiti:
\begin{enumerate}
    \item \textbf{Segretezza}: Le informazioni, nella loro precisa configurazione o in configurazioni dei loro elementi, non sono note o facilmente accessibili ad esperti e operatori di settore.
    \item \textbf{Valore economico}: Le informazioni in quanto segrete hanno valore economico: il mantenimento della loro segretezza conferisce un vantaggio economico a chi le detiene o viceversa la loro divulgazione causa perdite economiche
    \item \textbf{Misure di sicurezza adeguate}: le informazioni devono essere sottoposte a misure di sicurezza ragionevolmente adeguate, non solo dal punto di vista puramente tecnico ma anche organizzativo ie. prevedere corsi di formazione sulle misure a tuela del segreto.
\end{enumerate}

L'adeguatezza delle misure di sicurezza deve essere valutata in considerazione del contesto di applicazione, delle misre state-of-the-art e del valore del segreto che proteggono. è inoltre necessario che queste misure siano controllate e verificate periodicamente.

\begin{quote}
In contesti in cui per qualche motivo (principalmente economico) non è possibile istituire misure di sicurezza in un certo ambito (ie. protezione fisica) è necessario che si sopperisca in maniera alterativa.
\end{quote}

In contesto civilistico la misura di intervento giuridico è il \textbf{ricorso}: va presentato al giudice dimostrando (o dando forti elementi indiziari) la concorrenza sleale da parte di qualcuno soggetta all'utilizzo di segreti industriali. 
È necessario dimostrare che vi sia un'effettiva violazione del segreto e che le misure a sua protezione siano sufficienti.
Durante la fase preliminare del ricorso è sufficiente il \textit{fumus boni iuris}: un insieme di elementi che supportino sufficientemente le argomentazioni di chi lo ha esposto.
Va anche considerata la presenza di \textit{periculum in mora}: il mancato intervento causa un continuo/peggioramento dei danni a chi ha sporto ricorso.

A questo punto viene nominato un CTU e costituito un accesso \textit{inaudita altera  parte}: il CTU effettua la verifica senza preventiva notifica sotto ordine del giudice e il soggetto della verifica non può rifiutarsi. Questo processo si chiama \textit{descrizione}, termine che risale ad accertamenti di questo tipo in contesti puramente aziendali.

Se si dovesse riscontrare l'effettiva presenza delle informazioni segrete durante la descrizione si può procedere con il sequestro, l'applicazione di misure correttive specifiche (ritiro dal mercato e distruzione dei materiali derivati in modo abusivo) e la richiesta di risarcimento del danno. 

Nel caso del software, le procedure di descrizione, ispezione e sequestro seguono la Legge sul Diritto d'Autore anziché il CPI, e il CTU ha il compito di acquisire copie forensi non solo dell'eseguibile, ma anche di sorgenti, log e repository, documentando rigorosamente la catena di custodia.

\subsubsection{Quadro sanzionatorio penale}
In ambito penale la norma cardine è l'Art. 623 c.p. in materia alla rivelazione di segreti scientifici o industriali. A questa si affiancano i reati informatici strumentali, quali l'accesso abusivo a sistema informatico (Art. 615-ter c.p.), tipico vettore di esfiltrazione. Nel caso specifico del software, rileva anche l'Art. 171-bis L.d.A. per la duplicazione abusiva a fini di profitto.

\subsubsection{Diritto d'autore}
La legge a protezione del diritto d'autore nasce a tutela di opere proprietarie che non rimangono segrete (ie. opere d'arte, brani musicali, fotografie, ...).
In contesto informatico il codice sorgente può essere protetto cumulativamente come \textit{opera dell'ingegno} e, se mantenuto segreto, come \textit{segreto commerciale}.

Rientrano sotto questa protezione:  codice sorgente, eseguibili, documentazione progettuale.
Vengono protetti:
\begin{itemize}
    \item Programmi sorgente
    \item Programmi oggetto
    \item Materiale preparatorio (diagrammi di flusso, schemi logici, documentazione di progetto)
\end{itemize}

Non vengono protetti sotto il diritto d'autore:
\begin{itemize}
    \item Idee, principi ed algoritmi sottostanti
    \item Funzionalità del programma
    \item Linguaggi di programmazione e formati dei file
    \item Interfacce puramente funzionali
    \item Elementi obbligati dal contesto tecnico (\textit{scènes à faire}), legati indissolubilmente all'idea stessa (\textit{merger doctrine}) o appartenenti al pubblico dominio.
\end{itemize}

Il plagio del software può essere ``letterale'' (copia speculare del codice) o ``non letterale'' (copia della struttura, sequenza e organizzazione). Una tipica tecnica di sviluppo per prevenire accuse di plagio non letterale è il \textit{Clean Room Design}, che prevede l'uso di due team di sviluppo rigorosamente isolati (uno analizza l'originale redigendo le specifiche, il secondo programma senza mai vedere il codice sorgente di partenza).

Non è possibile brevettare software come stand-alone ma è possibile brevettarlo in combinazione con componenti hardware.

All'interno del diritto d'autore convivono due aspetti di tutela:
\begin{itemize}
    \item \textbf{Diritti Morali}: paternità, integrità, pubblicazione. Considerati inalienabili indipendentemente da accordi contrattuali.
    \item \textbf{Diritti Patrimoniali}: riproduzione, esecuzione, distribuzione, elaborazione. A differenza di quelli morali, sono cedibili.
\end{itemize}

La tutela del diritto d'autore ha scadenza di 70 anni, vengono tutelate anche rielaborazioni successive.
La tutela sul diritto d'autore non prevede la segretezza del codice ma, ove dimostrabile, protegge il lavoro originale dell'autore.

\subsubsection{Quantificazione del danno}
Tipologie di danno risarcibili:
\begin{itemize}
    \item Danno emergente: spese di indagine, sicurezza, limitazione del danno
    \item Lucro cessante: mancato guadagno attribuibile al periodo e all'oggetto
    \item Royalty equivalente: `prezzo del consenso' che il violatore avrebbe dovuto pagare (licenze)
    \item Danni morali
\end{itemize}

\paragraph{Contesto penale}
In diritto d'autore è applicabile anche in contesto penale in situazioni in cui a violazione viene utilizzata per generare profitto. 

\subsubsection{Titolarità del software}
In contesti di lavoro subordinato i diritti patrimoniali sul software spettano per legge al datore di lavoro (Art. 12-bis L.d.A.). Al contrario il diritto di paternità spetta sempre all'autore e non è cedibile in nessun contesto. Nel caso di software su commissione, la cessione dei diritti patrimoniali richiede necessariamente la forma scritta ad probationem. In contesti open source, i diritti possono essere in comunione tra più contributori.

\subsubsection{Brevetto}
Il brevetto tutela un informazione progettuale reso di pubblico dominio, proteggendo per un periodo limitato il \textit{lavoro d'ingegno} per 7 anni (so come stai facendo qualcosa ma non posso comunque replicarlo). 
Considerando la crescente integrazione del software in componenti meccaniche si è ampliata la protezione del brevetto anche ad esso ma esclusivamente come componente associato ad hardware. In Europa (tramite l'Ufficio Brevetti Europeo - EPO), le \textit{Computer-Implemented Inventions} (CII) sono brevettabili solo se presentano un contributo tecnico ulteriore (\textit{further technical effect}) che vada oltre la normale interazione fisica col computer.

\subsubsection{Modalità d'analisi (Software Forensics)}
Le metodologie d'esame si dividono in:
\begin{itemize}
    \item \textbf{Analisi Statica}: esame del codice senza esecuzione (sorgenti, script, metadati).
    \item \textbf{Analisi Dinamica}: osservazione a runtime (API, accessi al file system, traffico di rete).
    \item \textbf{Analisi Comparativa e Semantica}: confronto sistematico per individuare logiche identiche dissimulate.
\end{itemize}

Si fa un analisi delle risorse confrontando marcatori su più livelli:
\begin{itemize}
    \item Livello lessicale: string matching, commenti, nomi di variabile, costanti.
    \item Livello strutturale: confronto di alberi sintattici (AST diff) o Control Flow Graph (CFG) sui file binari.
    \item Storico di sviluppo (Git Forensics): analisi di commit, blame e repository, prestando attenzione a possibili manomissioni (rebase, alterazione dei timestamp).
\end{itemize}

Indizi forti di copia (le ``firme involontarie'') includono bug identici nella medesima posizione, errori di battitura nei commenti replicati e costanti errate preservate. Altri indicatori rilevanti provengono dalla Software Bill of Materials (SBOM) e dall'uso non dichiarato di librerie soggette a licenze Open Source. Il consulente tecnico (CT) incaricato ha il compito di descrivere le evidenze tecniche emerse e isolare il codice funzionalmente necessario, senza mai emettere autonomamente giudizi giuridici di ``plagio''.

\subsection{Reverse Engineering}
Il reverse engineering è vietato in quanto viola il diritto d'autore e può essere consentito dalla legge unicamente se strettamente necessario per conseguire l'interoperabilità con altri sistemi. 
L'utente legittimo può tuttavia analizzare il funzionamento black-box del software per comprenderne i principi, ma senza aggirare o eludere in maniera abusiva le misure poste a tutela del programma.

% --------------------------------------------------

\section{Valutazione Forense del Software}
La valutazione forense del software ha come obiettivo la determinazione del valore di mercato o del costo di riproduzione di un asset intangibile. A differenza di una stima ordinaria, la Digital Forensics (DF) impone un accertamento preliminare rigoroso sull'esistenza materiale, la consistenza tecnica, l'integrità e la titolarità del codice.

\subsection{Contesti di Applicazione e Inquadramento Civilistico}
La stima del software è rilevante in molteplici scenari:
\begin{itemize}
    \item \textbf{Fusioni e Acquisizioni (M\&A)}: per la due diligence tecnica.
    \item \textbf{Rivalutazione di asset}: obbligo di iscrizione tra le \textit{Immobilizzazioni Immateriali} seguendo principi contabili.
    \item \textbf{Stima del danno e Fallimentare}: quantificazione dei danni per plagio/pirateria o valorizzazione dell'asset da parte di curatori fallimentari (distinguendo tra valore di liquidazione e di continuità).
\end{itemize}

Un'errata valorizzazione espone al reato di falso in bilancio. In caso di conferimenti in S.p.A., è obbligatoria una \textit{perizia asseverata} nominata dal Tribunale.

\subsection{Fasi Preliminari e Certificazione Forense}
Prima di ogni stima economica, il perito deve certificare tre elementi fondamentali:
\begin{enumerate}
    \item \textbf{Autore}: identificazione di chi ha scritto il codice analizzando i log (Git, SVN).
    \item \textbf{Data Certa}: determinazione della marcatura temporale certificata di specifiche release.
    \item \textbf{Integrità}: verifica della corrispondenza sorgenti/binari tramite hash crittografico (SHA-256).
\end{enumerate}

\subsection{Standard e Metodologie di Stima}
La valutazione deve basarsi su standard internazionali, come \textbf{ISO/IEC 19770} (Software Asset Management) e \textbf{ISO/IEC 25010} (SQuaRE - per definire la qualità oggettiva del software in 8 caratteristiche).

\subsubsection{Criterio del Costo (Cost Approach)}
Valuta quanto costerebbe riprodurre il software da zero (ore/uomo, Function Points, costi di licenze e infrastrutture). 
\textit{Il ruolo della DF}: validare le ore dichiarate confrontandole con lo storico dei \textit{commit}, per smascherare ``gonfiamenti'' e distinguere codice nuovo da quello riciclato o copiato.

\subsubsection{Criterio di Mercato (Market Approach)}
Confronta l'asset con transazioni simili recenti calcolando multipli di valore (es. EV/Revenue). 
\textit{Il ruolo della DF}: fornire i parametri tecnici (stack, scalabilità, qualità) per giustificare gli aggiustamenti rispetto ai \textit{comparables}.

\subsubsection{Criterio Finanziario (Income Approach - DCF)}
Stima il valore tramite i flussi di cassa futuri attualizzati.
\textit{Il ruolo della DF}: valutare la qualità architetturale. Un codice disorganizzato (\textit{spaghetti code}) alza i costi di manutenzione e accresce il rischio di obsolescenza, riducendo drasticamente il valore dei flussi netti attesi.

\subsection{Analisi Tecnica e Metriche di Qualità}
Le metriche oggettive impattano direttamente sul valore economico:
\begin{itemize}
    \item \textbf{Metriche di Halstead}: misurano la complessità computazionale del codice.
    \item \textbf{Indice di Manutenibilità}: quantifica la facilità di intervento sul codice.
    \item \textbf{Debito Tecnico}: il costo implicito delle scorciatoie di sviluppo accumulate, considerato una \textit{passività latente}.
    \item \textbf{Metadati di sviluppo}: l'analisi di pipeline CI/CD o file di dipendenza rileva backdoor occulte, credenziali hard-coded e librerie obsolete.
\end{itemize}

\subsection{Fattori di Rischio e Svalutazione dell'Asset}
Il perito deve ponderare i seguenti elementi legali e tecnici che alterano il valore della stima:

\subsubsection{Licenze Open Source e SCA}
L'uso improprio di componenti ``copyleft'' (es. licenza GPL/AGPL) in software proprietario costringe al rilascio pubblico dei sorgenti derivati, potenzialmente azzerando il valore dell'asset. La prevenzione passa tramite una sistematica \textit{Software Composition Analysis (SCA)}.

\subsubsection{Cybersecurity e Privacy (GDPR)}
Vulnerabilità gravi (standard OWASP Top 10) rappresentano sconti diretti sul valore. L'assenza di principi \textit{Privacy by Design} espone a sanzioni amministrative fino al 4\% del fatturato, imponendo fondi rischi o pesanti svalutazioni strategiche.

\subsubsection{Proprietà Intellettuale (Diritti Morali vs Patrimoniali)}
La mancanza di chiari contratti di cessione dei diritti patrimoniali da parte di dipendenti o collaboratori freelance genera incertezza giuridica e conseguente deprezzamento della stima.

\subsection{Contesti Valutativi Particolari}
\begin{itemize}
    \item \textbf{Quantificazione dei danni}: somma di lucro cessante, danno emergente ed eventuale retroversione degli utili illeciti, dimostrati tramite log di installazioni illegittime e analisi di somiglianza.
    \item \textbf{Software Legacy}: la scarsità di professionisti per sistemi datati ma vitali (es. COBOL bancario) ne accresce enormemente il valore di continuità operativa.
    \item \textbf{Patent Box}: per l'accesso agli incentivi fiscali, la DF è necessaria per tracciare la ``lineage'' del codice e dimostrare il nesso causale tra costi di R\&S sostenuti e asset prodotto.
\end{itemize}

\subsection{Struttura della Perizia e Responsabilità}
La relazione tecnica si compone di quattro sezioni logiche: Identificazione Forense, Analisi Legale, Analisi Tecnica e Calcolo Economico.

Il perito estimatore assume responsabilità civile e penale. Vi è il preciso obbligo etico e legale di segnalare manipolazioni ad hoc del codice (codice ridondante, refactoring cosmetico pre-stima o falsi commit) orchestrate per incrementare artificialmente le metriche e le dimensioni del software.

\subsection{Frontiere Future}
\begin{itemize}
    \item \textbf{Intelligenza Artificiale}: la stima di un modello di AI impone metodologie doppie che separino il valore intrinseco del dataset di addestramento da quello dell'architettura generativa.
    \item \textbf{Blockchain}: l'uso di timestamping su blockchain (o NFT) diventerà lo standard per ottenere una prova di esistenza delle release software crittograficamente immutabile.
\end{itemize}

% --------------------------------------------------

\section{Gestione degli Incidenti e delle Prove Digitali: Standard ISO/IEC}
Questo documento sintetizza le linee guida operative e metodologiche fornite dagli standard internazionali ISO/IEC 27035 e ISO/IEC 27037, fondamentali per la gestione strutturata degli incidenti di sicurezza informatica e per la corretta acquisizione e conservazione delle prove digitali.

\subsection{Standard ISO/IEC 27035: Gestione degli Incidenti di Sicurezza}
L'ISO/IEC 27035 fornisce un approccio strutturato per identificare, valutare e rispondere agli incidenti informatici, con l'obiettivo di minimizzare l'impatto sul business e alimentare il miglioramento continuo del sistema ISMS (ISO/IEC 27001). 

La catena causale alla base dello standard prevede che una minaccia sfrutti una vulnerabilità di un asset, generando un \textit{evento} di sicurezza che, se impatta le operazioni, si trasforma in un \textit{incidente}.

\subsubsection{Ruoli Chiave}
\begin{itemize}
    \item \textbf{Point of Contact (PoC)}: Primo filtro aziendale che riceve le segnalazioni, effettua valutazioni preliminari e coordina l'intervento.
    \item \textbf{Incident Response Team (IRT)}: Team specialistico (noto anche come CERT o CSIRT) incaricato della gestione tecnica e operativa dell'incidente.
\end{itemize}

\subsubsection{Le 5 Fasi del Processo di Gestione}
Il processo è iterativo e si articola in cinque fasi strettamente interconnesse:
\begin{enumerate}
    \item \textbf{Pianificazione e preparazione}: Definizione di policy, costituzione dell'IRT, predisposizione di sistemi di logging e formazione del personale prima che l'incidente avvenga.
    \item \textbf{Individuazione e documentazione}: Monitoraggio continuo, raccolta informazioni e registrazione accurata degli eventi anomali, preservando le evidenze.
    \item \textbf{Valutazione e decisione}: Classificazione dell'evento per confermare se si tratti di un reale incidente e determinare il livello di risposta proporzionato (riducendo i falsi allarmi).
    \item \textbf{Risposta (Response)}: Fase operativa di contenimento, eliminazione della minaccia, recovery dei sistemi e chiusura dell'incidente. Se l'incidente non è sotto controllo, si innesca la gestione della crisi.
    \item \textbf{Apprendimento (Lessons Learned)}: Analisi post-incidente per identificare vulnerabilità, migliorare le procedure e aggiornare la valutazione del rischio.
\end{enumerate}

\subsubsection{Focus Operativo: Attacco Ransomware}
Il ransomware rappresenta l'incidente critico per eccellenza. Il ciclo di vita dell'attacco prevede: accesso iniziale (es. phishing), movimento laterale, esfiltrazione dati e infine la cifratura.
In caso di attacco:
\begin{itemize}
    \item \textbf{Contenimento}: Disconnessione immediata dalla rete dei dispositivi compromessi, valutando attentamente lo spegnimento per non perdere i dati volatili in RAM.
    \item \textbf{Recovery}: La presenza di un backup offline integro (strategia 3-2-1) è il fattore discriminante tra un ripristino rapido e la perdita totale/pagamento del riscatto.
    \item \textbf{Analisi}: Necessaria per verificare l'esfiltrazione e adempiere agli obblighi di notifica (Garante Privacy, ACN), analizzando artefatti di sistema (Event Log, MFT, Amcache).
\end{itemize}

\subsection{Standard ISO/IEC 27037: Gestione delle Prove Digitali}
L'ISO/IEC 27037 disciplina le procedure di primo intervento sulla scena del crimine informatico. Il suo perimetro è strettamente limitato all'identificazione, raccolta, acquisizione e conservazione della prova; \textbf{esclude deliberatamente la fase di analisi} per concentrarsi sulla protezione del dato da alterazioni.

\subsubsection{Ruoli Operativi}
\begin{itemize}
    \item \textbf{DEFR (Digital Evidence First Responder)}: Operatore qualificato che interviene per primo sulla scena, mettendo in sicurezza i dispositivi e preservando la catena di custodia.
    \item \textbf{DES (Digital Evidence Specialist)}: Esperto con competenze avanzate per gestire acquisizioni complesse (reti, sistemi cifrati, malware).
\end{itemize}

\subsubsection{Requisiti e Pilastri Metodologici}
Una prova digitale, per essere ammissibile, deve rispettare tre requisiti legali:
\begin{itemize}
    \item \textbf{Pertinenza}: Utilità diretta ai fini dell'indagine.
    \item \textbf{Affidabilità}: Prova genuina e non alterata.
    \item \textbf{Sufficienza}: Raccolta adeguata e proporzionata al caso.
\end{itemize}

Tutte le operazioni condotte dal DEFR/DES devono fondarsi su quattro pilastri: \textbf{Verificabilità, Ripetibilità, Riproducibilità e Giustificabilità} (ogni scelta tecnica deve poter essere difesa in sede di giudizio).

\subsubsection{Le 4 Fasi del Processo Forense}
\begin{enumerate}
    \item \textbf{Identificazione}: Riconoscimento dei dispositivi fisici (inclusi IoT, RFID, NAS) e dei dati volatili. La priorità assoluta è l'acquisizione dei dati in RAM.
    \item \textbf{Raccolta}: Rimozione e trasporto dei dispositivi. La regola d'oro in questa fase è: \textit{se un dispositivo è acceso, non va spento; se è spento, non va acceso}.
    \item \textbf{Acquisizione}: Creazione della \textit{copia forense certificata} (copia bit-a-bit che include spazio non allocato e slack space) apportando il minor numero possibile di alterazioni e documentandole rigorosamente.
    \item \textbf{Conservazione}: Protezione dei reperti da danni fisici (campi magnetici, scariche elettrostatiche).
\end{enumerate}

\subsubsection{La Catena di Custodia (Chain of Custody)}
Costituisce la storia completa del reperto. È un registro formale che traccia in modo ininterrotto chi, quando, come e perché ha avuto accesso al dispositivo dal momento della raccolta fino al processo. Include sempre i \textbf{valori di hash (MD5, SHA)} utilizzati come firma digitale per dimostrare matematicamente l'integrità dei dati originari. L'assenza o la compilazione errata della catena di custodia compromette irreparabilmente l'utilizzabilità della prova.

% --------------------------------------------------

\section{Digital Forensics Sanitaria}


L'obiettivo primario è la ricostruzione rigorosa di eventi clinici, amministrativi o tecnici, garantendo l'ammissibilità probatoria delle evidenze dinanzi all'autorità giudiziaria.

\begin{quote}
\textbf{NB.} L'attività forense in questo ambito si muove sull'intersezione di molteplici normative, in primis il GDPR, il D.Lgs. 101/2018, il Codice Deontologico Medico e il Regolamento MDR sui dispositivi medici.
\end{quote}

Per le indagini su reati gravi, dove vige un conflitto tra il segreto professionale medico e l'obbligo di referto, l'acquisizione dei sistemi clinici necessita sempre di un decreto motivato dall'autorità giudiziaria e del rispetto del principio di proporzionalità.

\subsection{Aree di indagine e sistemi critici}
I bersagli dell'indagine forense spaziano dai complessi software gestionali dell'ente ospedaliero ai singoli dispositivi a stretto contatto con il paziente.

\subsubsection{Cartella Clinica Elettronica (CCE)}
La CCE costituisce un atto pubblico di fede privilegiata: l'alterazione fraudolenta di questi documenti configura il reato di falso in atto pubblico. 
Qualsiasi analisi presuppone:
\begin{itemize}
    \item \textbf{Controllo delle Firme}: l'obbligo normativo di utilizzare la Firma Elettronica Qualificata (FEQ) o Avanzata (FEA) per garantire l'immodificabilità del referto dopo la sottoscrizione.
    \item \textbf{Analisi dell'Audit Trail}: l'investigazione sistematica dei log di sistema per tracciare rigorosamente ogni accesso, modifica o cancellazione.
    \item \textbf{Verifica delle Inconsistenze Temporali}: il confronto tra i timestamp dell'applicazione, del DBMS e del sistema operativo per rilevare modifiche retroattive non visibili dalla normale interfaccia utente.
\end{itemize}

\subsubsection{Dispositivi Medici (IoMT) e Diagnostica}
L'estrazione di dati da dispositivi loMT (es. pacemaker, pompe insuliniche, robot chirurgici) mira a verificare se un evento avverso sia scaturito da un malfunzionamento, una manomissione o un errore clinico.
\begin{itemize}
    \item \textbf{Formato DICOM}: i ricchi metadati presenti nelle immagini diagnostiche consentono di svelare steganografie forensi, manipolazioni di radiografie o discordanze sulle date dell'esame.
    \item \textbf{Software Forensics}: indagini basate sull'analisi di bug o sulla revisione dei log di input/output dei sensori per verificare i comandi impartiti a macchinari autonomi.
\end{itemize}

\begin{quote}
L'accesso fisico per l'estrazione di dati da dispositivi medici impiantati nel paziente richiede un'autorizzazione giudiziaria specifica e la supervisione di un medico specialista.
\end{quote}

\subsection{Reati e Responsabilità}
La ricostruzione forense affianca e integra la classica perizia medico-legale.

\subsubsection{Controlli difensivi e Cyber-Responsabilità}
Accertamenti eseguiti per individuare condotte illecite dei professionisti o negligenze della struttura sanitaria stessa.
\begin{itemize}
    \item \textbf{Accesso abusivo a sistema informatico}: la prassi illegittima di consultare i dati clinici di colleghi o figure pubbliche senza un'effettiva necessità medica (Art. 615-ter c.p.). L'illiceità si dimostra incrociando i log di accesso con i registri degli appuntamenti.
    \item \textbf{Cyber-Responsabilità dell'ente}: in scenari di attacchi ransomware o data breach, viene valutata la colpa organizzativa (D.Lgs. 231/2001) e la presenza documentata di adeguati piani di Business Continuity o misure di sicurezza tecniche a tutela del paziente.
    \item \textbf{Furto di identità medica}: la sottrazione di un template biometrico (usato per firmare o accedere a medicinali) costituisce un danno permanente e irreversibile in quanto la biometria, a differenza di una password, non può essere ``reimpostata''.
\end{itemize}

\subsection{Esonero dalle procedure standard}
A causa delle specificità sanitarie, il normale modus operandi delle indagini forensi subisce delle radicali variazioni.

\subsubsection{Territorio grigio e Incident Response}
La necessità di non spegnere macchine salvavita obbliga gli investigatori in ambito ospedaliero a far ricorso esclusivo alla ``Live Forensics''. Si procede all'acquisizione dei log e della memoria volatile in tempo reale, aggirando le classiche procedure di contenimento per rispettare il vincolo etico primario delle cure vitali.

\begin{quote}
In contesti affini, strumenti personali non prettamente sanitari (es. fitness tracker e smartwatch) producono tracce utilissime in ambito di medicina legale digitale per stimare con accuratezza l'esatto momento del decesso o la caduta dei parametri vitali.
\end{quote}

% --------------------------------------------------

\section{ Domande d'esame}

\begin{itemize}
    \item \textbf{Superamento della prova di resistenza giudiziaria} 
    \begin{itemize}
        \item L'onere della prova dell'attendibilità del dato è interamente a carico di chi lo produce in giudizio, a causa della natura intrinsecamente manipolabile dei bit streams.
    \end{itemize}
    
    \item \textbf{Differenza procedimento/processo penale} 
    \begin{itemize}
        \item Il procedimento penale è la fase delle indagini preliminari (dominata dal PM), il processo penale è eventuale ed è la fase che inizia solo con l'esercizio dell'azione penale (davanti al Giudice).
    \end{itemize}
    
    \item \textbf{Differenza perizia (penale) e CTU (civile)} 
    \begin{itemize}
        \item Nel penale la perizia è un mezzo di prova autonomo (ammesso solo per questioni complesse), mentre nel civile la CTU è uno strumento integrativo e ausiliario per colmare le lacune tecniche del giudice.
    \end{itemize}
    
    \item \textbf{Accertamenti ripetibili vs irripetibili} 
    \begin{itemize}
        \item I secondi modificano il reperto irreversibilmente (es. sblocco distruttivo, analisi di un sistema live) e impongono l'obbligo di avviso alla difesa per garantire il contraddittorio fin da subito.
    \end{itemize}
    
    \item \textbf{La catena di custodia e l'hash} 
    \begin{itemize}
        \item L'acquisizione forense deve garantire l'inalterabilità tramite calcolo di impronte crittografiche (es. SHA-256) e documentazione ininterrotta di ogni passaggio.
    \end{itemize}
    
    \item \textbf{Come la marca temporale garantisce i file di log al momento dell'evento} 
    \begin{itemize}
        \item Firma crittografica da parte di ente esterno che garantisce l'hash in un preciso momento, anche a certificato scaduto (LTV).
    \end{itemize}
    
    \item \textbf{Tipologie di acquisizione di un disco rigido} 
    \begin{itemize}
        \item \textbf{Fisica (bit-stream image)}: copia esatta bit a bit dell'intero supporto, inclusi spazi non allocati e file cancellati; \textbf{Logica}: estrazione dei soli file allocati e visibili al file system; \textbf{Selettiva/Mirata}: raccolta di specifici artefatti (utile in fase di triage); \textbf{Live}: acquisizione a sistema acceso, necessaria per catturare memoria RAM e connessioni di rete attive.
    \end{itemize}
    
    \item \textbf{Necessità dell'applicazione della marca temporale anche ai log di acquisizione} 
    \begin{itemize}
        \item Oltre all'immagine forense, i file di log certificano \textit{come} e \textit{con quali esiti} è avvenuta l'acquisizione tecnica. La marca temporale su questi log ne ``cristallizza'' l'integrità e la datazione in modo opponibile a terzi, impedendo contestazioni legali su eventuali alterazioni successive o retrodatazioni del verbale digitale.
    \end{itemize}
    
    \item \textbf{Quando viene nominato il CTP} 
    \begin{itemize}
        \item Nel \textbf{civile}, le parti lo nominano tipicamente entro un termine assegnato dal giudice dopo la nomina del CTU (Consulente Tecnico d'Ufficio). Nel \textbf{penale}, può essere nominato in ogni fase (purché non superi il numero dei periti/consulenti del PM), ma è fondamentale nominarlo \textbf{prima} dell'inizio di accertamenti tecnici irripetibili (es. art. 360 c.p.p.) per poter partecipare e garantire il contraddittorio durante le operazioni distruttive.
    \end{itemize}
    
    \item \textbf{Analisi dei duplicati nelle firme digitali} 
    \begin{itemize}
        \item In ambito forense si riferisce principalmente all'uso degli algoritmi di hash (le ``firme'') per la \textbf{deduplicazione}: confrontare enormi quantità di file per escludere quelli di sistema noti (es. tramite database NSRL) o identificare copie esatte dello stesso file. Include anche l'analisi delle ``file signature'' (magic numbers) per smascherare file rinominati con estensioni false per occultare informazioni.
    \end{itemize}
    
    \item \textbf{La procedura di descrizione in un processo civile quando viene rubato un segreto commerciale} 
    \begin{itemize}
        \item Si ricorre tipicamente alla \textbf{``descrizione inaudita altera parte''} (senza avvisare la controparte) ex art. 129/130 del Codice della Proprietà Industriale o 692 c.p.c. Il giudice autorizza l'Ufficiale Giudiziario, affiancato da un CTU informatico, a fare un'ispezione a sorpresa (es. nella sede del sospettato) per estrarre copie forensi o dati mirati \textit{prima} che i dispositivi possano essere formattati o le prove distrutte.
    \end{itemize}
    
    \item \textbf{Definzione e disegno dello slackspace} 
    \begin{itemize}
        \item Lo \textit{slack space} (spazio di risulta) è lo spazio residuo inutilizzato che si crea tra la fine logica di un file e la fine fisica del cluster (l'unità minima di allocazione sul disco) che lo ospita. Poiché il sistema operativo considera l'intero cluster occupato, questo spazio non viene sovrascritto e può conservare frammenti di dati di file precedenti (cancellati), assumendo enorme rilevanza forense per recuperare password o vecchi documenti.
    \end{itemize}
    
    \item \textbf{Definzione e descrizioni delle iso (35 -> gestione degli incidenti) e (37 -> gestione delle prove digitali)}
    \begin{itemize}
        \item \textbf{ISO/IEC 27035}: Standard per l'Incident Management della sicurezza informatica. Fornisce un approccio strutturato in 5 fasi (preparazione, identificazione, valutazione, risposta e \textit{lessons learned}) per rilevare e mitigare attacchi e violazioni in ambito aziendale. 
        \item \textbf{ISO/IEC 27037}: Linee guida rigorose destinate ai \textit{first responder} per l'identificazione, raccolta, acquisizione e conservazione della Digital Evidence. Detta i principi metodologici affinché la prova digitale mantenga il suo valore probatorio e sia legalmente ammissibile in tribunale.
    \end{itemize}
\end{itemize}

\end{document}
